\section{One Hida ideal and the complete selected Count-II support input}
\label{sec:hida-installation-v45}

This chapter assembles the original objects, after their producers
have been reconstructed.  It has two different scopes.  For every
original positive rank it constructs the whole
\(\mathrm{HidaUAlt}\): one ideal with all six Hida properties
and the entire same-ideal Taylor--Wiles automorphy system.  For
each of the two rank-four Links used in Count II it also supplies
the original \(\mathrm{P1core}\) on that ideal.  The latter uses
the completed one-point support proof.  The general assertion
\(\mathrm{PfromAut}\), with arbitrary chains of
\(\mathrm{NiceE}\) points and its \(\mathrm{P3core}\) conclusion,
is a separate antecedent in the all-rank modularity implication
below.  It is not inferred from the one-point argument.

Retain the original \(H'=\mathrm{HypothesesH}'\), including its
named automorphic, local-correspondence and cohomological
foundations.  In particular the original definite-inner-form
base change, cuspidal Galois attachment and printed local
classification inputs retain their previous status.  The
cohomological package \(\mathrm{CohP}\) means all eight original
fields: local finiteness, local Tate duality, local Euler
characteristic, global finiteness, global Euler characteristic,
Greenberg--Wiles, normalized Poitou--Tate, and the original tame
local theorem.  The full \(\mathrm{LocCFT}\) and the actual
Artin normalizations are retained, with their already reconstructed
finite consumers.  No general support or automorphy assertion is
added to these foundations.

For the all-rank statements fix the original
\(\ell,L,\Omega,s,V,c,n,r,B,\Lambda,\iota\), where
\(n\geq1\), \(O=c.O\), \(V.S'=s.S\), and
\(\iota:\overline{\mathbb Q}_\ell\simeq\mathbb C\) is the
given field isomorphism.  All coefficient maps, residual labels,
decomposition groups and local character orders are the original
ones.  There is no continuity requirement on \(\iota\).
The original setting already has \(\ell\ne2\); the prime three
remains included throughout.

% BEGIN HIA1
\paragraph{HIA1. All six Hida fields on the same given ideal.}
\label{claim:HIA1}
Let \(R\) be any of the original unitary Hecke data, let \(K\)
be its given auxiliary level, and suppose
\(R.\mathrm{IsJ}(B,K,J)\).  Then, under the retained original
inputs and \(V.S'=s.S\), the literal record
\[
                         B.\mathrm{HidaPropsAlt}(\iota,J,\Lambda)
                         \tag{HIA.1}
\]
holds for this very ideal \(J\).
\begin{proof}
Use the following six proved endpoints with the same original
arguments and the same ideal:
\[
\begin{array}{c|c|l}
\text{field}&\text{proved endpoint}&\text{literal content}\\ \hline
\mathrm{ring}&\mathrm{JRP15}&\mathrm{RingProps}\ \Lambda\ J\\
\mathrm{dense}&\mathrm{UPA18}&\mathrm{Dense}\ J\\
\mathrm{pt}&\mathrm{HPT18}&\mathrm{PointTheoremBaseObject}\ \iota\ J\\
\mathrm{end1}&\mathrm{HED16}&\mathrm{ArithmeticPointDecomposition}\ \iota\ J\\
\mathrm{end2}&\mathrm{HAW19}&\mathrm{IrrArithPointsAutomorphicAway}\ \iota\ J\\
\mathrm{twist}&\mathrm{HTS8}&\mathrm{TwistOk}\ J.
\end{array}
\tag{HIA.2}
\]
The original aliases \(\mathrm{arithmeticPointPurityDense}\)
and \(\mathrm{arithmeticPointClassicalityDense}\) are exactly
the fourth and fifth contents of this table.  No implication
from a weaker purity or classicality predicate is being used.

There are several substantive guards inside these fields.
The point theorem quantifies over a supplied representation
and proves \(J\subseteq\ker x\) under its original conditions
\(2\leq n\leq4\), continuity, automorphy, full auxiliary
inertia triviality, and the prescribed ordered flags and
nonlinkage.  It is an implication with those rank guards even
when the ambient rank is outside that interval.  The first
end field retains every positive summand rank and positive
string length, their exact total rank, the full group
characteristic-polynomial identity, and its strict cofinite
root-ratio assertion over all original places under the original
condition that every string length is one.  The second
retains \(\mathrm{AutCsd0}\) without a continuity antecedent;
only its \(\mathrm{WitOff}\) conclusion is conditional on
continuity.  Its one real window center is chosen before all
places for each attached cusp.  The twist field is the
original finite \(\ell\)-power-order twist assertion, on this same ideal.

The ring field retains the original \(r.\mathrm{V0ok}\)
antecedent in its auxiliary-unramified clause; that antecedent
is not imposed on the entire six-field record.  HPT18 and
HAW19 likewise add no such hypothesis to their endpoints.
Every endpoint in (HIA.2) accepts the given
\(\mathrm{IsJ}\), rather than selecting an unrelated ideal.
Thus the six proofs are the six fields of (HIA.1).

This also includes the legitimate zero-module branch
\(J=B\).  Fields whose antecedents contain
\(J\subseteq\ker x\) then have no such point.  The
\(\mathrm{pt}\) field does not have that antecedent and is
not declared vacuous: its full conclusion is precisely the
one supplied by HPT18 for this \(\mathrm{IsJ}\) branch.
\end{proof}
% END HIA1

% BEGIN HIA2
\paragraph{HIA2. The full all-rank Hida and automorphy supplier.}
\label{claim:HIA2}
Under the original \(H'\), the entire \(\mathrm{HidaUAlt}\)
holds.  In particular, for every original datum in its definition,
including \(V.S'=s.S\), \(\mathrm{noMu}\),
\(\mathrm{Large}(L,c)\), every \(n\geq1\), and every
actual \(r,B,\Lambda,\iota\), there is one ideal \(J\subseteq B\)
such that
\[
 B.\mathrm{HidaPropsAlt}(\iota,J,\Lambda)
                   \quad\hbox{and}\quad \mathrm{HasAut}(B,J).
 \tag{HIA.3}
\]
No further coefficient extension is required.
\begin{proof}
Take \(Q=\mathrm{QData.empty}(V,S'=S)\).  ITA13 gives
the actual eight-field unitary Hecke datum \(R\) over the
specified \(O\) and \(\Lambda\).  There is a unique tuple
of auxiliary levels indexed by its empty set; use it as \(K\).
Put
\[
 H=\mathrm{hset}(Q,R.\tau_0),\qquad
 M=H.\mathrm{MK}(R.\theta,R.a,K),\qquad
 \mathbb T=H.\mathrm{TK}(R.\theta,R.a,K).
 \tag{HIA.4}
\]
HLM9 supplies the literal \(\mathrm{IsJ}\) on these objects
for every rank in the statement.  If \(M=0\), it takes
\(J=B\).  Otherwise it constructs the continuous original
determinant and local characters on \(\mathbb T\), proves
their full pair relations and their prescribed weight map,
and obtains a surjection \(f:B\to\mathbb T\).  It takes
\(J=\ker f\).  Since \(\mathbb T\) is its actual subalgebra
of \(\operatorname{End}_{\Lambda}(M)\), this is also
\(\operatorname{Ann}_B(M)\).  Thus the support ideal has been
chosen before any later arithmetic point or auxiliary
Taylor--Wiles family.  HIA1 gives all six fields on it.

JRP1 verifies \(\mathrm{IsCoef}(\ell,O)\): the given
coefficient ring is a complete discrete valuation ring with
finite residue field of characteristic \(\ell\).
The original local package used in converting levels to
automorphy is exactly
\[
 \mathrm{LocPAlt}=
 (\text{Jacquet finite length/admissibility},
   \text{geometric lemma},\text{tame local theorem}).
 \tag{HIA.5}
\]
Its first two inputs are the original literature fields
\(H'.\mathrm{litN.loc.bz}\) and
\(H'.\mathrm{litN.loc.geo}\).  Its third is the original
tame-local consequence of the retained local Tate duality
field.  These are the same named foundations as in HTW24;
none is inferred merely from the six Hida fields.

HLM26 constructs the full \(\mathrm{LevelAlt}(B,J)\) with
one base \(b\) before every auxiliary datum.  HTW24 converts
that very family, with (HIA.5), into the whole
\(\mathrm{HasAut}(B,J)\).  In expanded form this says
\[
 \begin{gathered}
 \forall\mathcal T:\mathrm{TWGrp}(\mathrm{globOf}(r,V)),
 \quad\mathrm{TWPl}(r,\mathcal T)
 \ \Longrightarrow\\
 \mathrm{Nonempty}\!\left(
 \mathrm{TWAut}(\mathrm{pairOf}(B.u),J,\mathrm{Empty},
                   n[L^+:\mathbb Q],\mathcal T)\right).
 \end{gathered}
 \tag{HIA.6}
\]
In each system the residual eigenvalue lifts are chosen before
all its levels; the exponent is the original \(n!\), uniform
in the level, and both original trace-polynomial relations
and the coefficient-map correction are retained.  No
division by \(n!\) occurs.  The zero-module branch is the
one already covered by the full level and automorphy proofs.

The choices above use arbitrary supplied \(B\) and
\(\Lambda\); universality is not restricted to one preferred
representative.  They take place over the original \(O\),
and therefore work for every coefficient object in the
definition, including every sufficiently large later object.
Together with (HIA.1), (HIA.6) proves exactly (HIA.3) and
the full quantifier order of \(\mathrm{HidaUAlt}\).
\end{proof}
% END HIA2

% BEGIN HIA3
\paragraph{HIA3. The cyclotomic condition and the separate modularity implication.}
\label{claim:HIA3}
For every original level \(V\), its \(\mathrm{noMu}\)
condition implies the actual residual cyclotomic condition
\[
 \exists g\in V.\mathrm{frame}.G_0,
                 \overline\varepsilon(g)\ne1.
 \tag{HIA.7}
\]
Independently of the selected rank-four support proof below,
the original all-rank implication is
\[
 \mathrm{CohP}+\mathrm{LocCFTAll}+\mathrm{PfromAut}
              +\mathrm{HidaUAlt}
                       \ \Longrightarrow\ \mathrm{ModUAlt}.
 \tag{HIA.8}
\]
Here \(\mathrm{PfromAut}\) is an explicit antecedent of
(HIA.8), not a conclusion of this chapter.
\begin{proof}
Choose a primitive \(\ell\)th root \(\zeta\) in the
given algebraic closure.  If every element of \(G_L\)
fixed \(\zeta\), the fixed-field theorem would put it
in \(L\), contrary to \(\mathrm{noMu}\).  Thus some
\(\sigma\in G_L\) moves it.  The action on these roots
is multiplication by the reduction of the cyclotomic
character, so \(\overline\varepsilon(\sigma)\ne1\).
By the defining cyclotomic compatibility of the original
level projection, the element \(g=V.\mathrm{pr}(\sigma)\)
lies in the original \(G_0\) and has that same reduction.
This proves (HIA.7).  It does not construct a new level
quotient or replace its cyclotomic character.

To prove (HIA.8), fix any input in the literal definition
of \(\mathrm{ModUAlt}\).  The \(\mathrm{HidaUAlt}\)
antecedent supplies one \(J\) satisfying (HIA.3); HIA2
has now supplied that antecedent under the original \(H'\).
Use JRP1 for \(\mathrm{IsCoef}\), and use
\(\mathrm{LocCFTAll}\) on this \(s\).
The given \(n\geq1\), \(V.S'=s.S\), (HIA.7), the full
\(\mathrm{CohP}\), and \(\mathrm{HasAut}(B,J)\) are
precisely the antecedents of the original
\(\mathrm{PfromAut}\).  That hypothesis returns
\[
                       B.\mathrm{P1core}(J)
                          \quad\hbox{and}\quad
                       B.\mathrm{P3core}(J).
 \tag{HIA.9}
\]
Appending these two fields to \(\mathrm{HidaPropsAlt}\)
is exactly \(B.\mathrm{ModularAlt}(\iota,J,\Lambda)\).
This proves (HIA.8) at all ranks and on the same ideal.

For clarity about its boundary, \(\mathrm{P3core}\)
quantifies over every finite chain of \(\mathrm{NiceE}\)
points with the original thick connecting primes.  The
one-point \(\mathrm{NiceI}\) argument that follows has
no such connecting primes and supplies only its stated
\(\mathrm{P1core}\) conclusion.  It does not discharge
the \(\mathrm{PfromAut}\) antecedent of (HIA.8).
\end{proof}
% END HIA3

% BEGIN HIA4
\paragraph{HIA4. The selected rank-four support conclusion on that ideal.}
\label{claim:HIA4}
Fix either original rank-four Link, its actual \(B,\Lambda\),
and the ideal \(J\) obtained in HIA2.  Retain the full
original \(\mathrm{CohP}\) and local reciprocity hypotheses.
Then the entire original one-point predicate
\(B.\mathrm{P1core}(J)\) holds:
\[
 \begin{gathered}
 x.\mathrm{NiceI},\quad B.\mathrm{IsPt}(x,f),\quad
 J\subseteq\ker f,\quad
 \mathfrak q\in\operatorname{Spec}B,\quad
 \mathfrak q\subseteq\ker f\\
                       \Longrightarrow\quad J\subseteq\mathfrak q.
 \end{gathered}
 \tag{HIA.10}
\]
All the maps, primes and point fields in (HIA.10) are the
original ones.  This conclusion requires no general
\(\mathrm{PfromAut}\) or \(\mathrm{P3core}\).
\begin{proof}
We verify the completed inputs to JSP11 in their actual order.
JRP1 supplies the coefficient condition.  KX3 supplies the
whole \(\mathrm{LocCFT}\) on the chosen decomposition
groups; CIP12 supplies the exact additional Artin matching
and the already reconstructed normalization consumer.
The Link has \(V.S'=s.S\) and \(\mathrm{noMu}\), so the
unconditional first assertion of HIA3 supplies (HIA.7).
The required \(\mathrm{CohP}\) is the retained eight-field
package, and HIA2 supplies \(\mathrm{HasAut}(B,J)\)
before \(x,f,\mathfrak q\) are introduced.

Fix arbitrary \(x,f,\mathfrak q\) as in (HIA.10).
CIP10 and CIP11 normalize this \(\mathrm{NiceI}\) point
on the actual updated inertia frame.  They give the original
\(a=4\), \(N3\), and \(N4(D_0)\), where \(D_0\)
is the original auxiliary subgroup, together with an
injective local map of point rings and the full determinant
and ordered-character identities.  The first part of CAP13
supplies the required \(\Phi\)-finiteness before a finite
list is selected.  CCH15 then supplies that list, its second
point extension and its common Kap with \(\kappa.g=0\).
CCH13 preserves the original \(N4(D_0)\), the complete
determinant law and the ordered characters.  If \(j\)
denotes the composite point-ring map, it is injective, has
the original compatible coefficient structure, and gives
\[
                             g=jf,\qquad\ker g=\ker f.
 \tag{HIA.11}
\]
Only these pointwise finite extensions are made; the outer
coefficient object and its earlier threshold stay fixed.

Now use the second part of CAP13 to realize precisely the
selected list by the actual \(\mathcal T\) and
\(\mathrm{TWPl}\).  The residual-kernel identity of this
record makes every residual character one on its subgroup
\(H\).  The full determinant identity therefore gives
\(\operatorname{cp}(\overline{\rho(h)})=(X-1)^4\)
for \(h\in H\).  JSP11's Cayley--Hamilton and
characteristic-\(\ell\) argument applies on this same
point: \((\rho(h)-1)^4\in tM_4(k[[t]])\), whence
\(\rho(h)^{\ell^{4m}}-1\in t^mM_4(k[[t]])\).
Taking the open normal cores and then every \(m\)
proves the exact \(\mathrm{PL}(\ell,H)\)-kernel
condition.  Thus the later bounds receive the full pro-\(\ell\)
condition, not only a residual statement.

CFO20 and CFO21 now give the required first-order field
on this very point and Kap, with parameter
\[
                            4[L^+:\mathbb Q]+4\mathcal T.q.
 \tag{HIA.12}
\]
Its constant is chosen after the point and list, before
the truncation and eventual level.  CCH16 gives the
obstruction field on the same Kap with constant zero;
CCH17 checks their coupled installation.  No independent
Kap or alternative point is used for the two bounds.
Together with the original eigenvalue and detector data
these construct the literal \(\mathrm{Pt}\) required
by the support consumer.

JSP11 supplies the actual \(\mathrm{Adapted}\),
\(\mathrm{GenBy}\), stable-discriminant and kernel
comparisons.  Its use of \(\mathrm{HasAut}\) is now
fed by HIA2 at this \(\mathcal T/\mathrm{TWPl}\), with
the same \(J\), the empty parameter type and
dimension \(4[L^+:\mathbb Q]\).  SCI13 supplies the
full selected \(\mathrm{Raw}\) conclusion for exactly
\(h=0,a=n=4,\kappa.g=0\), including its actual
quotient, coefficient normalization and prime contraction.
Apply JSP11 with these supplied inputs and initial stable
prime \(\ker g\).  Its conclusion is
\(J\subseteq\mathfrak q\), after (HIA.11) identifies
the original point kernel.  This proves (HIA.10) for
every original test.

All the hypotheses formerly left in the selected route
have now been supplied: normalization, choice, actual
auxiliary primes, the coupled first- and second-order
bounds, selected Raw, and same-ideal HasAut.  The arithmetic
and cohomological hypotheses printed above are unchanged.
The proof does not invert \(24\), exclude \(\ell=3\),
or assert the support theorem for arbitrary chains.
\end{proof}
% END HIA4

% BEGIN HIA5
\paragraph{HIA5. Two excluded places give one global window.}
\label{claim:HIA5}
Let \(n\geq1\), let \(u\ne u'\) be finite places of the
maximal real subfield, and use the same \(\iota\) and
the same representation \(\rho\).  Then the exact
original implication is
\[
 \begin{split}
 &\mathrm{AutOffN}(u,\iota,\rho)\ \wedge\
       \mathrm{AutOffN}(u',\iota,\rho)\\
 &\hspace{12mm}\Longrightarrow\
 \mathrm{AutCsd0}(\iota,\rho)\ \wedge\
       \bigl(\mathrm{Continuous}(\rho)\Longrightarrow
                                  \mathrm{WitA}(\iota,\rho)\bigr).
 \end{split}
 \tag{HIA.13}
\]
In particular its global window has one real center before
all eligible finite places; the two input centers need not agree.
\begin{proof}
Keep the \(\mathrm{AutCsd0}\) part of the first input.
Only for the remaining conditional assertion assume
\(\rho\) continuous.  The two inputs give
\(\mathrm{WitOff}(u,\iota,\rho)\) and
\(\mathrm{WitOff}(u',\iota,\rho)\).  Either gives the
first part of \(\mathrm{WitA}\): at every place above
\(\ell\), \(\mathrm{Unl1}\) implies
\(\mathrm{PotUnl}\), with all original prolongations
and positive powers.

Fix any regular algebraic cusp \(\pi\) attached to
this \(\rho\).  HAW6 supplies a single real \(t\),
before any place is considered, such that toral unlinked
support at a place implies \(\mathrm{WinAt}(\pi,t,v)\).
The first and second \(\mathrm{WitOff}\) also supply
centers \(a\) and \(b\) for their respective excluded-place
conditions.  Take an arbitrary eligible place \(v\), meaning
\(v\nmid\ell\) or \(\mathrm{PotUnl}(\rho,v)\).
If its underlying finite place in the maximal real subfield
is different from \(u\), use the first window; otherwise
that place is different from \(u'\),
so use the second.  In either case the resulting window
gives toral support and, for every support character,
\[
                    q_v^{c-1/2}<|\xi(\varpi_v)|<q_v^{c+1/2}
                     \quad(c=a\text{ or }b).
 \tag{HIA.14}
\]
If two such characters had unramified quotient taking the
value \(q_v\) on \(\varpi_v\), their absolute values
would differ by the factor \(q_v\).  The smaller strict
lower bound in (HIA.14) would then put the larger one
strictly above its upper bound.  Thus the support is
unlinked in the exact original sense.  Repeated equal
characters cause no difficulty.  Apply the already chosen
\(t\) from HAW6 to get \(\mathrm{WinAt}(\pi,t,v)\).
This works for every eligible \(v\) without a new choice
of \(t\).  It proves the second part of \(\mathrm{WitA}\)
for every attached cusp, and hence (HIA.13).
\end{proof}
% END HIA5

% BEGIN HIA6
\paragraph{HIA6. Both original Links above one common coefficient threshold.}
\label{claim:HIA6}
Retain the original Fit/Deep tower, the original named
foundations, and a fixed \(\iota\).  The entire selected
\(\mathrm{COUNT\_II\_SUPPORT}\) input is now supplied.
More precisely, there is one original coefficient threshold
\(c_1\), chosen before every later \(c'\), Space,
universal ring and weight object, such that both original
rank-four Links have the following property at every
allowed \(c'\): for every actual \(B,\Lambda\) there
is one \(J\subseteq B\) carrying all six Hida fields
and the full \(\mathrm{P1core}\).  The Links retain
their common \(\kappa,\beta,t\), distinct auxiliary
primes, and every original local and residual field.
\begin{proof}
Fix the tower \(D,T\) and \(\iota\) first.  QX4 constructs
its characters before choosing an auxiliary prime, and
chooses \(c_1\) containing their required finite coefficient
values and all the original embeddings of \(\mathcal O_L\).
It follows that for every \(c'\) containing both the original
\(c\) and \(c_1\), the literal Links
\(\mathcal L_0,\mathcal L_1\) exist over \(c'\), with
the same \(\kappa,\beta,t\) and with distinct auxiliary
rational primes.  This coefficient choice precedes all
Spaces and all their universal data.

Each Link includes exactly the hypotheses needed here:
\(V.S'=s.S\), \(\mathrm{noMu}\),
\(\mathrm{Large}(L,c')\), the actual residual rank-four
object, and \(r.\mathrm{V0ok}\), as well as the original
local-place identification and character equalities.
For either Link and arbitrary supplied \(B,\Lambda\),
apply HIA2 over that same \(c'.O\).  It gives a single
ideal with all six fields and full HasAut.  HIA4 adds
\(\mathrm{P1core}\) to that exact ideal.  The cohomological
input is the retained full \(\mathrm{CohP}\); the local
class-field and cyclotomic conditions have been checked
in HIA4.  The type, ordinary and arithmetic-point
constructions introduce no coefficient enlargement.
The normalization and support extensions occur only inside
the later \(\mathrm{P1core}\) point quantifier.  Hence no
new external coefficient threshold is needed beyond QX4's.

The resulting quantifier order is therefore
\[
 \begin{gathered}
 D,T,\iota\ ;\quad \exists c_1\ ;\quad
 \forall c'\supset c,c_1\ ;\quad
 \exists\mathcal L_0,\mathcal L_1\ ;\\
 \forall i\in\{0,1\}\ \forall B_i,\Lambda_i\ ;\quad
 \exists J_i\ ;\quad
 B_i.\mathrm{HidaPropsAlt}(\iota,J_i,\Lambda_i)
                    \ \wedge\ B_i.\mathrm{P1core}(J_i).
 \end{gathered}
 \tag{HIA.15}
\]
It remains valid for every later original Space.  Each
\(J_i\) is chosen before the point, prime, detector,
truncation and level tests in its properties.  The two
ideals belong to their two respective original universal
rings; no equality between ideals in different rings is
asserted.  Within each ring all properties have one witness.
This is exactly the same-ideal and two-Link contract isolated
by JSP12, now with its previously retained producers supplied.

To exhibit its original consumer, take an arbitrary original
Space \(X\) above this threshold.  QX5 constructs the
actual nonempty ordinary pairs, so KX9 supplies actual
universal rings, KX5 supplies their weight objects, and
QX5 gives the maps \(\phi_i:B_i\to X.B\) for \(a=0\).
Use (HIA.15) and set
\[
                    I=\phi_0(J_0)X.B+\phi_1(J_1)X.B.
 \tag{HIA.16}
\]
QX10 requires precisely the six Hida fields and selected
\(\mathrm{P1core}\) now proved, rather than an arbitrary-chain
\(\mathrm{P3core}\).  It proves all five original fields
\(\mathrm{fin},\mathrm{p1},\mathrm{seed},\mathrm{pur},\mathrm{aut}\)
of \(\mathrm{PotentialModularityTower}\) on this same
\(X,I\), and thus the entire \(\mathrm{PMExAlt}\).
The already proved YX2 implication then gives the original
\(\mathrm{CountII}\) under its unchanged background inputs.

The stronger end2 fields also preserve the original two-place
gluing step explicitly.  Whenever the source consumer gives
\(\mathrm{AutOffN}\) on both Links, QX5's comparison and
their common \(\kappa\) identify its input as the same
twisted representation.  Distinct auxiliary rational primes
give distinct auxiliary finite places of the maximal real subfield:
equality of the latter
would make their absolute norms, and hence their degree-one
rational primes, equal.  HIA5 therefore gives the full
global window with one center.  This observation neither
identifies unrelated representations nor requires the two
original window centers to coincide.

All these conclusions retain \(\ell=3\), the fixed
\(\iota\), full determinant laws and pro-\(\ell\) conditions,
and the original cohomological and literature boundaries.
They close this selected support producer.  They do not
prove the general \(\mathrm{PfromAut}\) antecedent in
HIA3 or remove the other independent obligations in the
complete RH and GRH reconstruction.
\end{proof}
% END HIA6
